\documentclass[12pt,a4paper]{article}

% ─── Packages ───────────────────────────────────────────────────────────────
\usepackage[T1]{fontenc}
\usepackage[utf8]{inputenc}
\usepackage[bahasa]{babel}
\usepackage{geometry}
\geometry{margin=2.5cm}

\usepackage{setspace}
\onehalfspacing

\usepackage{parskip}
\setlength{\parindent}{1.25cm}
\setlength{\parskip}{6pt}

\usepackage{titlesec}
\titleformat{\section}{\large\bfseries}{\thesection.}{0.5em}{}
\titleformat{\subsection}{\normalsize\bfseries}{\thesubsection.}{0.5em}{}
\titleformat{\subsubsection}{\normalsize\itshape}{\thesubsubsection.}{0.5em}{}

\usepackage{listings}
\usepackage{xcolor}

\definecolor{codebg}{rgb}{0.97,0.97,0.97}
\definecolor{javakw}{rgb}{0.13,0.13,0.67}
\definecolor{javastr}{rgb}{0.20,0.55,0.20}
\definecolor{javacom}{rgb}{0.45,0.45,0.45}
\definecolor{javaann}{rgb}{0.55,0.20,0.60}

\lstdefinestyle{javaStyle}{
  backgroundcolor=\color{codebg},
  basicstyle=\ttfamily\small,
  keywordstyle=\color{javakw}\bfseries,
  stringstyle=\color{javastr},
  commentstyle=\color{javacom}\itshape,
  emph={@Test,@ParameterizedTest,@ValueSource,@MethodSource,@CsvFileSource,
        @BeforeAll,@BeforeEach,@AfterEach,@AfterAll},
  emphstyle=\color{javaann}\bfseries,
  breaklines=true,
  frame=single,
  framerule=0.4pt,
  rulecolor=\color{gray!50},
  tabsize=4,
  showstringspaces=false,
  captionpos=b,
  numbers=left,
  numberstyle=\tiny\color{gray},
  numbersep=8pt,
  xleftmargin=1.5em,
  framexleftmargin=1.5em,
}
\lstset{style=javaStyle}

\usepackage{booktabs}
\usepackage{array}
\usepackage{longtable}
\usepackage{hyperref}
\hypersetup{colorlinks=true, linkcolor=black, urlcolor=blue, citecolor=black}

\usepackage{caption}
\captionsetup[lstlisting]{font=small, labelfont=bf}

% ─── Metadata ────────────────────────────────────────────────────────────────
\title{%
  \textbf{Analisis Cara Kerja dan Fungsi Kode\\
  pada \textit{Parameterized Testing} dengan JUnit~5}\\[4pt]
  \large Studi Kasus: Proyek \texttt{Unit\_Testing1} --- Wallet Management System
}
\author{%
  Pemrograman Perangkat Lunak Lanjut (PPPL)\\
  Tugas 1 --- Pembahasan Kode
}
\date{2025}

% ─── Document ────────────────────────────────────────────────────────────────
\begin{document}

\maketitle
\tableofcontents
\newpage

% =============================================================================
\section{Pendahuluan}
% =============================================================================

Pengujian perangkat lunak merupakan salah satu tahapan kritis dalam siklus pengembangan
perangkat lunak.  Tanpa pengujian yang memadai, kebenaran fungsional suatu sistem tidak
dapat dijamin, sehingga potensi kegagalan di lingkungan produksi menjadi jauh lebih
besar.  Dalam konteks \textit{unit testing}, setiap unit fungsional yang terkecil---biasanya
sebuah metode---diverifikasi secara terisolasi guna memastikan perilakunya sesuai dengan
spesifikasi yang telah ditetapkan.

Repositori \texttt{PPPL\_Tugas1} ini mengimplementasikan sistem manajemen dompet
(\textit{Wallet Management System}) sederhana yang ditulis dalam bahasa Java, beserta
rangkaian pengujian unit yang komprehensif menggunakan kerangka kerja
JUnit~5.\footnote{JUnit~5 adalah versi terkini dari kerangka kerja pengujian unit Java
yang terdiri atas tiga modul utama: JUnit Platform, JUnit Jupiter, dan JUnit Vintage.}
Pengujian dalam proyek ini secara khusus memanfaatkan fitur \textit{parameterized testing},
yaitu kemampuan untuk menjalankan satu definisi metode uji terhadap sejumlah nilai
masukan yang berbeda-beda secara otomatis, tanpa duplikasi kode.

Dokumen ini membahas secara akademis cara kerja dan fungsi setiap komponen kode dalam
repositori tersebut, dengan penekanan pada mekanisme \textit{parameterized testing} yang
diterapkan, meliputi anotasi \texttt{@ValueSource}, \texttt{@MethodSource}, dan
\texttt{@CsvFileSource}.

% =============================================================================
\section{Struktur Proyek}
% =============================================================================

Proyek dikonfigurasi menggunakan Apache Maven dengan Java~21 sebagai versi target
kompilasi.  Dependensi utama yang dideklarasikan dalam berkas \texttt{pom.xml} adalah
\texttt{junit-jupiter} versi 5.10.2, yang menyediakan seluruh API pengujian yang
digunakan.  Secara keseluruhan, struktur kode terbagi menjadi dua lapisan: kode produksi
(\textit{main}) dan kode pengujian (\textit{test}).

\begin{center}
\begin{tabular}{ll}
\toprule
\textbf{Lapisan} & \textbf{Berkas} \\
\midrule
Kode produksi & \texttt{Wallet.java}, \texttt{Owner.java}, \texttt{InsufficientFundsException.java} \\
Kode pengujian & \texttt{WalletTest.java} \\
Data uji & \texttt{withdraw-valid.csv}, \texttt{withdraw-invalid.csv} \\
\bottomrule
\end{tabular}
\end{center}

% =============================================================================
\section{Kelas-Kelas Produksi}
% =============================================================================

\subsection{Kelas \texttt{InsufficientFundsException}}

Kelas ini merupakan \textit{custom exception} yang mewarisi kelas
\texttt{RuntimeException}.  Keberadaannya mencerminkan prinsip desain yang baik: alih-alih
menggunakan pengecualian generik, sistem mendefinisikan tipe pengecualian khusus yang
secara semantik merepresentasikan kondisi saldo tidak mencukupi.  Konstruktor tunggalnya
menerima sebuah pesan deskriptif yang kemudian diteruskan ke konstruktor induk melalui
pemanggilan \texttt{super(message)}.

\begin{lstlisting}[language=Java, caption={Definisi \texttt{InsufficientFundsException}}]
public class InsufficientFundsException extends RuntimeException {
    public InsufficientFundsException(String message) {
        super(message);
    }
}
\end{lstlisting}

Karena merupakan \textit{unchecked exception} (turunan \texttt{RuntimeException}), kelas
ini tidak memaksa pemanggil untuk menanganinya secara eksplisit, sehingga kode klien
tetap bersih dari blok \texttt{try-catch} yang tidak diperlukan.

\subsection{Kelas \texttt{Owner}}

Kelas \texttt{Owner} merepresentasikan pemilik dompet dengan dua atribut tidak dapat
diubah (\textit{immutable}): \texttt{name} dan \texttt{email}.  Kedua atribut ini
dideklarasikan sebagai \texttt{final} dan hanya dapat ditetapkan melalui konstruktor,
yang memastikan bahwa setiap objek \texttt{Owner} yang telah dibuat bersifat konsisten
sepanjang masa hidupnya.

\begin{lstlisting}[language=Java, caption={Definisi \texttt{Owner}}]
public class Owner {
    private final String name;
    private final String email;

    public Owner(String name, String email) {
        this.name = name;
        this.email = email;
    }

    public String getName()  { return name; }
    public String getEmail() { return email; }
}
\end{lstlisting}

Desain \textit{immutable} ini sangat relevan dalam konteks pengujian karena objek
\texttt{Owner} dapat digunakan kembali sebagai data masukan lintas pengujian tanpa
risiko perubahan keadaan yang tidak diinginkan.

\subsection{Kelas \texttt{Wallet}}

Kelas \texttt{Wallet} adalah kelas utama yang mengenkapsulasi logika bisnis sistem
dompet.  Kelas ini mengelola tiga aspek utama: informasi pemilik (\textit{owner}),
koleksi kartu (\textit{cards}), dan daftar uang tunai (\textit{cashList}).

\subsubsection{Manajemen Pemilik}

Kelas ini mendukung dua varian penetapan pemilik: penetapan melalui nama string biasa
dan penetapan melalui objek \texttt{Owner}.  Ketika pemilik ditetapkan melalui objek
\texttt{Owner}, atribut \texttt{ownerName} juga ikut diperbarui secara otomatis dari
\texttt{owner.getName()}, sehingga kedua representasi selalu sinkron.

\subsubsection{Manajemen Kartu}

Metode \texttt{addCard(String card)} hanya menerima kartu yang tidak bernilai
\texttt{null} dan tidak kosong, sebagaimana tercermin dari kondisi validasi berikut:

\begin{lstlisting}[language=Java, caption={Validasi pada \texttt{addCard}}]
public void addCard(String card) {
    if (card != null && !card.isEmpty()) {
        cards.add(card);
    }
}
\end{lstlisting}

Metode \texttt{takeCard(String card)} menghapus dan mengembalikan kartu yang sesuai
dari koleksi; jika kartu tidak ditemukan, metode mengembalikan \texttt{null}.  Metode
\texttt{getCards()} mengembalikan salinan defensif dari daftar kartu
(\texttt{new ArrayList<>(cards)}) sehingga perubahan pada daftar yang dikembalikan tidak
memengaruhi keadaan internal dompet.

\subsubsection{Manajemen Uang Tunai}

Metode \texttt{addCash(int amount)} hanya menambahkan nilai yang positif.  Metode
\texttt{takeCash(int amount)} menghapus entri pertama yang cocok secara nilai; karena
penggunaan \texttt{cashList.remove(Integer.valueOf(amount))}, metode ini menghapus
berdasarkan nilai (\textit{by value}), bukan berdasarkan indeks.

Metode \texttt{withdrawCash(int amount)} mengimplementasikan logika penarikan yang lebih
kompleks.  Ia terlebih dahulu memvalidasi bahwa jumlah penarikan positif (melempar
\texttt{IllegalArgumentException} bila tidak), kemudian memeriksa kecukupan saldo total
(melempar \texttt{InsufficientFundsException} bila saldo kurang).  Proses penarikan
mengiterasikan \texttt{cashList} secara berurutan, mengumpulkan entri yang akan dihapus
hingga jumlah target terpenuhi, dan mengembalikan kembalian (\textit{change}) bila ada.

\begin{lstlisting}[language=Java, caption={Logika \texttt{withdrawCash}}]
public void withdrawCash(int amount) {
    if (amount <= 0)
        throw new IllegalArgumentException("Jumlah withdraw harus lebih dari 0");
    int total = getTotalCash();
    if (amount > total)
        throw new InsufficientFundsException("Saldo tidak cukup");
    int remaining = amount;
    List<Integer> toRemove = new ArrayList<>();
    for (int cash : cashList) {
        if (remaining <= 0) break;
        toRemove.add(cash);
        remaining -= cash;
    }
    cashList.removeAll(toRemove);
    int change = -remaining;
    if (change > 0) cashList.add(change);
}
\end{lstlisting}

% =============================================================================
\section{Kelas Pengujian: \texttt{WalletTest}}
% =============================================================================

\subsection{Gambaran Umum}

\texttt{WalletTest} adalah kelas pengujian unit yang menggunakan API JUnit~5 Jupiter.
Kelas ini tidak bersifat \texttt{public} (sesuai konvensi JUnit~5 yang tidak
memerlukannya), dan mendeklarasikan atribut instansi \texttt{wallet} bertipe
\texttt{Wallet} yang selalu diinisialisasi ulang sebelum setiap metode uji dijalankan.

\subsection{Siklus Hidup Pengujian (\textit{Test Lifecycle})}

JUnit~5 menyediakan empat anotasi siklus hidup yang digunakan dalam kelas ini:

\begin{center}
\begin{tabular}{lll}
\toprule
\textbf{Anotasi} & \textbf{Cakupan} & \textbf{Fungsi} \\
\midrule
\texttt{@BeforeAll}  & Sekali per kelas & Mencetak pesan pembuka sesi pengujian \\
\texttt{@BeforeEach} & Sebelum tiap uji & Membuat instansi \texttt{Wallet} baru \\
\texttt{@AfterEach}  & Setelah tiap uji & Mengatur ulang \texttt{wallet} ke \texttt{null} \\
\texttt{@AfterAll}   & Sekali per kelas & Mencetak pesan penutup sesi pengujian \\
\bottomrule
\end{tabular}
\end{center}

Pola ini sangat penting untuk menjamin \textit{test isolation}: setiap metode uji
mendapatkan objek \texttt{Wallet} yang bersih, sehingga keadaan yang ditinggalkan oleh
satu pengujian tidak memengaruhi pengujian berikutnya.  Hal ini merupakan praktik standar
dalam \textit{unit testing} yang mencegah terjadinya \textit{test coupling}.

\begin{lstlisting}[language=Java, caption={Anotasi siklus hidup pada \texttt{WalletTest}}]
@BeforeAll
static void initAll() {
    System.out.println("=== Mulai test WalletTest ===");
}

@BeforeEach
void setUp() {
    wallet = new Wallet();
}

@AfterEach
void tearDown() {
    wallet = null;
}

@AfterAll
static void tearDownAll() {
    System.out.println("=== Semua test WalletTest selesai ===");
}
\end{lstlisting}

Perlu dicatat bahwa metode yang dianotasi \texttt{@BeforeAll} dan \texttt{@AfterAll}
harus bersifat \texttt{static} karena dipanggil pada level kelas, bukan pada level
instansi.

% =============================================================================
\section{Parameterized Testing}
% =============================================================================

\textit{Parameterized testing} adalah teknik penulisan pengujian di mana satu metode uji
dijalankan berulang kali dengan nilai masukan yang berbeda-beda, yang disuplai oleh sebuah
\textit{provider}.  Dalam JUnit~5, teknik ini diaktifkan dengan anotasi
\texttt{@ParameterizedTest} yang dikombinasikan dengan salah satu anotasi sumber data
(\textit{source annotation}).  Pendekatan ini secara signifikan mengurangi duplikasi kode
dibandingkan dengan penulisan metode uji terpisah untuk setiap nilai masukan.

Proyek ini menggunakan tiga jenis sumber data parameterized testing, yaitu
\texttt{@ValueSource}, \texttt{@MethodSource}, dan \texttt{@CsvFileSource}.  Masing-masing
memiliki karakteristik, kelebihan, dan konteks penggunaan yang berbeda, sebagaimana
diuraikan pada subbagian berikut.

\subsection{\texttt{@ValueSource}: Sumber Nilai Tunggal}

\texttt{@ValueSource} merupakan cara paling sederhana untuk menyuplai data ke metode
parameterized testing.  Anotasi ini mendukung array dari tipe primitif maupun
\texttt{String} dan \texttt{Class}, dan setiap elemen array tersebut akan digunakan
sebagai argumen tunggal untuk satu eksekusi metode uji.

\subsubsection{Pengujian Nilai Uang Tunai Valid}

\begin{lstlisting}[language=Java, caption={Pengujian \texttt{addCash} dengan nilai valid menggunakan \texttt{@ValueSource}}]
@ParameterizedTest
@ValueSource(ints = {1000, 5000, 10000, 50000, 100000})
void testAddCashValid(int amount) {
    wallet.addCash(amount);
    assertEquals(1, wallet.getCashList().size());
    assertTrue(wallet.getCashList().contains(amount));
}
\end{lstlisting}

Metode ini dieksekusi sebanyak lima kali, sekali untuk setiap nilai dalam array
\texttt{\{1000, 5000, 10000, 50000, 100000\}}.  Pada setiap eksekusi, JUnit~5 menyuplai
nilai yang bersangkutan sebagai argumen \texttt{amount}.  Karena \texttt{@BeforeEach}
membuat objek \texttt{Wallet} baru sebelum setiap eksekusi, setiap iterasi dimulai
dengan dompet yang kosong.

Setiap eksekusi memverifikasi dua hal: pertama, bahwa ukuran \texttt{cashList} bertambah
menjadi tepat satu (artinya uang berhasil ditambahkan); kedua, bahwa nilai tersebut
memang terdapat dalam daftar.  Kelima nilai yang diuji merupakan denominasi uang kertas
Rupiah yang umum, menunjukkan desain uji yang relevan secara bisnis.

\subsubsection{Pengujian Nilai Uang Tunai Tidak Valid}

\begin{lstlisting}[language=Java, caption={Pengujian \texttt{addCash} dengan nilai tidak valid menggunakan \texttt{@ValueSource}}]
@ParameterizedTest
@ValueSource(ints = {0, -1000, -50000})
void testAddCashInvalid(int amount) {
    wallet.addCash(amount);
    assertEquals(0, wallet.getCashList().size(),
        "Cash negatif/nol tidak boleh masuk");
}
\end{lstlisting}

Metode ini menguji tiga nilai batas: nol, nilai negatif kecil, dan nilai negatif besar.
Masing-masing diharapkan \emph{tidak} mengubah ukuran \texttt{cashList}, karena logika
validasi dalam \texttt{Wallet.addCash()} menolak nilai yang tidak lebih besar dari nol
(\texttt{if (amount > 0)}).  Pesan kustom pada argumen ketiga \texttt{assertEquals}
memperjelas ekspektasi apabila pengujian gagal.

\subsubsection{Ringkasan Mekanisme \texttt{@ValueSource}}

\texttt{@ValueSource} sangat efektif untuk menguji satu parameter primitif dengan
sejumlah kecil nilai yang dapat didaftar secara \textit{inline}.  Kelemahannya adalah
tidak dapat digunakan untuk menguji lebih dari satu parameter sekaligus, sehingga untuk
skenario multiparameter diperlukan anotasi sumber data lainnya.

\subsection{\texttt{@MethodSource}: Sumber Data dari Metode Pabrik}

\texttt{@MethodSource} memungkinkan metode uji menerima parameter kompleks, termasuk
objek Java, dengan cara menunjuk ke sebuah metode \textit{factory} yang mengembalikan
\texttt{Stream<Arguments>} atau struktur sejenis.

\subsubsection{Pengujian Penetapan Objek \texttt{Owner}}

\begin{lstlisting}[language=Java, caption={Metode pabrik data untuk \texttt{@MethodSource}}]
static Stream<Arguments> ownerProvider() {
    return Stream.of(
        Arguments.of(new Owner("Budi",  "budi@mail.com")),
        Arguments.of(new Owner("Sari",  "sari@mail.com")),
        Arguments.of(new Owner("Andi",  "andi@mail.com"))
    );
}
\end{lstlisting}

\begin{lstlisting}[language=Java, caption={Pengujian \texttt{setOwner(Owner)} dengan \texttt{@MethodSource}}]
@ParameterizedTest
@MethodSource("ownerProvider")
void testSetOwnerObject(Owner owner) {
    wallet.setOwner(owner);

    assertNotNull(wallet.getOwnerObject(),
        "Owner object tidak boleh null");
    assertSame(owner, wallet.getOwnerObject());
    assertEquals(owner.getName(), wallet.getOwner());
    assertEquals(owner.getEmail(),
        wallet.getOwnerObject().getEmail());
}
\end{lstlisting}

Metode \texttt{ownerProvider()} adalah metode statis yang mengembalikan aliran tiga
objek \texttt{Arguments}, masing-masing membungkus sebuah instansi \texttt{Owner} yang
berbeda.  JUnit~5 mengiterasi aliran ini dan memanggil \texttt{testSetOwnerObject}
sebanyak tiga kali, sekali untuk setiap objek \texttt{Owner}.

Dalam setiap eksekusi, empat asersi diverifikasi:
\begin{enumerate}
  \item \textbf{\texttt{assertNotNull}} --- memastikan bahwa setelah \texttt{setOwner},
        referensi objek \texttt{Owner} dalam dompet tidak \texttt{null}.
  \item \textbf{\texttt{assertSame}} --- memverifikasi identitas referensi objek, yaitu
        bahwa \texttt{wallet.getOwnerObject()} mengembalikan \emph{referensi yang sama
        persis} dengan objek \texttt{owner} yang disuplai, bukan salinan baru.  Ini
        membuktikan bahwa \texttt{setOwner(Owner)} menyimpan referensi asli, bukan
        melakukan kloning.
  \item \textbf{\texttt{assertEquals}} pertama --- memverifikasi bahwa
        \texttt{ownerName} dalam dompet sinkron dengan nama dari objek \texttt{Owner}
        yang baru ditetapkan.
  \item \textbf{\texttt{assertEquals}} kedua --- memverifikasi bahwa data email dapat
        diakses melalui jalur \texttt{getOwnerObject().getEmail()}.
\end{enumerate}

\subsubsection{Kapan Menggunakan \texttt{@MethodSource}}

\texttt{@MethodSource} adalah pilihan yang tepat ketika: (a) parameter pengujian berupa
objek Java (bukan tipe primitif atau \texttt{String}); (b) data pengujian perlu dibuat
secara programatis; atau (c) terdapat lebih dari satu parameter dengan tipe yang beragam.
Metode pabrik harus bersifat \texttt{static} dan harus dideklarasikan dalam kelas yang
sama atau dapat direferensikan secara penuh.

\subsection{\texttt{@CsvFileSource}: Sumber Data dari Berkas CSV Eksternal}

\texttt{@CsvFileSource} memungkinkan data pengujian dipisahkan sepenuhnya dari kode
pengujian ke dalam berkas CSV yang disimpan di direktori sumber daya (\textit{resources}).
Pendekatan ini sangat menguntungkan ketika data pengujian berjumlah besar, perlu dikelola
oleh pemangku kepentingan non-teknis, atau perlu diperbarui tanpa mengubah kode sumber.

\subsubsection{Pengujian Penarikan Uang Valid}

Data uji untuk skenario valid disimpan dalam berkas
\texttt{src/test/resources/withdraw-valid.csv}:

\begin{lstlisting}[caption={Isi berkas \texttt{withdraw-valid.csv}}]
deposit,withdraw,expectedTotal
10000,0,10000
20000,5000,15000
50000,25000,25000
100000,100000,0
\end{lstlisting}

Baris pertama berupa baris tajuk (\textit{header}) yang dilewati berkat parameter
\texttt{numLinesToSkip = 1}.  Empat baris berikutnya mendefinisikan empat skenario
pengujian: tanpa penarikan, penarikan sebagian, penarikan setengah, dan penarikan penuh.

\begin{lstlisting}[language=Java, caption={Pengujian penarikan valid dengan \texttt{@CsvFileSource}}]
@ParameterizedTest
@CsvFileSource(resources = "/withdraw-valid.csv",
               numLinesToSkip = 1)
void testWithdrawCashValid(int deposit,
                           int withdraw,
                           int expectedTotal) {
    wallet.addCash(deposit);
    if (withdraw > 0) {
        wallet.withdrawCash(withdraw);
    }
    assertEquals(expectedTotal, wallet.getTotalCash(),
        "Deposit " + deposit
        + " - withdraw " + withdraw
        + " harus = " + expectedTotal);
}
\end{lstlisting}

Untuk setiap baris CSV, JUnit~5 menyuplai tiga kolom sebagai tiga argumen metode.
Kondisi \texttt{if (withdraw > 0)} mencegah pemanggilan \texttt{withdrawCash} dengan
nilai nol (yang akan melempar \texttt{IllegalArgumentException}), karena baris pertama
data memang bermaksud menguji skenario tanpa penarikan.  Pesan kustom pada
\texttt{assertEquals} mempermudah diagnosa kegagalan dengan menampilkan konteks data
saat itu.

\subsubsection{Pengujian Penarikan Uang Tidak Valid}

Data uji untuk skenario tidak valid disimpan dalam berkas
\texttt{src/test/resources/withdraw-invalid.csv}:

\begin{lstlisting}[caption={Isi berkas \texttt{withdraw-invalid.csv}}]
initialDeposit,withdraw,exceptionType
10000,15000,InsufficientFundsException
10000,0,IllegalArgumentException
10000,-1000,IllegalArgumentException
0,1000,InsufficientFundsException
\end{lstlisting}

Setiap baris mendefinisikan satu skenario kesalahan yang diharapkan: (1) penarikan
melebihi saldo; (2) penarikan dengan nilai nol; (3) penarikan dengan nilai negatif; dan
(4) penarikan dari dompet tanpa dana.

\begin{lstlisting}[language=Java, caption={Pengujian penarikan tidak valid dengan \texttt{@CsvFileSource}}]
@ParameterizedTest
@CsvFileSource(resources = "/withdraw-invalid.csv",
               numLinesToSkip = 1)
void testWithdrawCashInvalid(int initialDeposit,
                             int withdraw,
                             String exceptionType) {
    if (initialDeposit > 0) {
        wallet.addCash(initialDeposit);
    }
    if (exceptionType.equals("InsufficientFundsException")) {
        assertThrows(InsufficientFundsException.class,
            () -> wallet.withdrawCash(withdraw),
            "Harus throw InsufficientFundsException");
    } else {
        assertThrows(IllegalArgumentException.class,
            () -> wallet.withdrawCash(withdraw),
            "Harus throw IllegalArgumentException");
    }
}
\end{lstlisting}

Metode ini membaca kolom ketiga (\texttt{exceptionType}) sebagai \texttt{String} dan
menggunakannya untuk menentukan tipe pengecualian yang diharapkan secara dinamis.
\texttt{assertThrows} dari JUnit~5 memverifikasi bahwa ekspresi lambda yang diberikan
melempar pengecualian dari tipe yang tepat; jika tidak ada pengecualian yang dilempar
atau tipenya berbeda, pengujian dinyatakan gagal.

Penggunaan kolom \texttt{exceptionType} berbasis string ini merupakan pola yang umum
dalam parameterized testing berbasis CSV untuk menghindari keterbatasan bahwa berkas CSV
tidak dapat merepresentasikan tipe Java secara langsung.

\subsubsection{Keunggulan \texttt{@CsvFileSource}}

Pemisahan data dari kode memberikan beberapa keuntungan: (a) data pengujian dapat
dikembangkan secara independen dari kode; (b) mempermudah tinjauan kasus uji oleh analis
bisnis; dan (c) memungkinkan penambahan skenario baru tanpa kompilasi ulang kode
pengujian---hanya berkas CSV yang perlu diperbarui.

% =============================================================================
\section{Perbandingan Mekanisme \textit{Parameterized Testing}}
% =============================================================================

Tabel berikut merangkum perbedaan karakteristik ketiga mekanisme parameterized testing
yang digunakan dalam proyek ini:

\begin{longtable}{>{\raggedright\arraybackslash}p{3cm}
                  >{\raggedright\arraybackslash}p{3cm}
                  >{\raggedright\arraybackslash}p{4cm}
                  >{\raggedright\arraybackslash}p{4cm}}
\toprule
\textbf{Anotasi} & \textbf{Tipe Parameter} & \textbf{Kelebihan} & \textbf{Keterbatasan} \\
\midrule
\endhead
\texttt{@ValueSource}    & Primitif / \texttt{String}  &
  Sintaks ringkas, mudah dibaca \textit{inline} &
  Hanya satu parameter, tidak mendukung objek \\[4pt]
\texttt{@MethodSource}   & Objek Java arbitrer &
  Fleksibel, mendukung konstruksi programatis &
  Memerlukan metode pabrik statis tambahan \\[4pt]
\texttt{@CsvFileSource}  & Primitif / \texttt{String} (dari CSV) &
  Data terpisah dari kode, mudah dikelola &
  Terbatas pada tipe yang dapat direpresentasikan CSV \\
\bottomrule
\end{longtable}

% =============================================================================
\section{Pengujian Konvensional (\textit{Non-Parameterized})}
% =============================================================================

Di samping pengujian terparameterisasi, \texttt{WalletTest} juga memuat sejumlah metode
uji konvensional yang menggunakan anotasi \texttt{@Test} biasa.  Pengujian-pengujian ini
menangani skenario yang bersifat tunggal dan tidak memerlukan variasi nilai masukan,
seperti verifikasi bahwa \texttt{owner} bernilai \texttt{null} saat dompet baru dibuat,
atau bahwa \texttt{getCards()} mengembalikan salinan defensif dan bukan referensi langsung
ke koleksi internal.

Secara khusus, metode \texttt{testGetCardsReturnSalinan} (Listing~\ref{lst:defensive})
memverifikasi perilaku salinan defensif:

\begin{lstlisting}[language=Java,
                   caption={Verifikasi salinan defensif pada \texttt{getCards()}},
                   label={lst:defensive}]
@Test
void testGetCardsReturnSalinan() {
    wallet.addCard("KTP");
    wallet.getCards().add("SIM");
    assertEquals(1, wallet.getCards().size(),
        "getCards() harus return salinan");
}
\end{lstlisting}

Pengujian ini menambahkan sebuah kartu ke dompet, kemudian mencoba menambahkan kartu
lain melalui daftar yang dikembalikan oleh \texttt{getCards()}.  Ekspektasi bahwa ukuran
daftar tetap satu membuktikan bahwa modifikasi pada daftar kembalian tidak memengaruhi
keadaan internal dompet.

Adapun metode \texttt{testWalletLengkap} merupakan pengujian integrasi ringan yang
memverifikasi kombinasi operasi: penetapan pemilik, penambahan kartu, dan penambahan uang
tunai dalam satu skenario terpadu.

% =============================================================================
\section{Kesimpulan}
% =============================================================================

Repositori \texttt{PPPL\_Tugas1} mendemonstrasikan penerapan \textit{parameterized
testing} secara menyeluruh menggunakan JUnit~5 Jupiter.  Melalui tiga mekanisme sumber
data---\texttt{@ValueSource}, \texttt{@MethodSource}, dan \texttt{@CsvFileSource}---kode
pengujian berhasil menghindari redundansi sekaligus mencakup berbagai skenario masukan
secara sistematis.

Anotasi \texttt{@ValueSource} terbukti efisien untuk validasi nilai primitif tunggal
dengan cakupan nilai batas yang representatif.  Anotasi \texttt{@MethodSource}
memungkinkan pengujian terhadap objek Java kompleks dengan jaminan identitas referensi
yang diverifikasi melalui \texttt{assertSame}.  Sementara itu, \texttt{@CsvFileSource}
mewujudkan pemisahan yang tegas antara kode dan data uji, yang mendukung prinsip
\textit{separation of concerns} dan memudahkan pemeliharaan jangka panjang.

Penggunaan siklus hidup pengujian yang tepat---\texttt{@BeforeEach} untuk inisialisasi
dan \texttt{@AfterEach} untuk pembersihan---memastikan isolasi antaruji sehingga setiap
kasus uji bersifat mandiri dan hasilnya dapat dipercaya.  Secara keseluruhan, praktik
pengujian dalam proyek ini mencerminkan penerapan prinsip-prinsip \textit{clean testing}
yang sejalan dengan standar industri pengembangan perangkat lunak modern.

\end{document}
